\section{Related work}
\label{sec:related}

\paragraph{Random convolutional kernels.} ROCKET classifies time series with a linear
model on features pooled from thousands of random convolutional kernels
\citep{dempster2020rocket}. MiniRocket replaces the random kernels by a small fixed set with
two weight values \citep{dempster2021minirocket}, MultiRocket adds pooling operators and the
first difference \citep{tan2022multirocket}, and Hydra organises random kernels into
competing groups \citep{dempster2023hydra}. QUANT, an interval method, classifies quantiles
over fixed intervals with extremely randomised trees, in less than 15 minutes of
single-core compute for 142 UCR datasets \citep{dempster2024quant}. In the
most recent bake-off, the concatenation of MultiRocket and Hydra features and the HIVE-COTE
2.0 ensemble were the most accurate methods on the UCR archive
\citep{middlehurst2024bakeoff,middlehurst2021hc2}. Two lines of work reduce the cost of
these transforms after the fact: pruning kernels or features, which can remove most of them
with little loss of accuracy \citep{salehinejad2022srocket,uribarri2024detach}, and
combining fast transforms with stronger models, such as stacked ensembles of Hydra and
QUANT \citep{maniar2025metalearning} or tabular in-context learners on ROCKET features
\citep{orourke2026rocketpfn}. \hh{} takes a different route: it changes what a group is, so
that the cost of finding a group's winner no longer grows with its size. Its features could
enter the same combinations.

\paragraph{Low-precision and edge variants.} The closest work is NanoHydra
\citep{cioflan2025nanohydra}, which runs Hydra on a microcontroller with binary kernels
($\pm1$), integer arithmetic and a quantised linear classifier. Its groups remain random
sets of kernels that are evaluated one by one. In its own evaluation on the UCR archive,
NanoHydra averages 83.81\% against 85.31\% for Hydra (Hydra's number as reported by
\citet{schaefer2023weasel2}), although its numbers are the best of several runs per
dataset. The two approaches address different costs: NanoHydra makes each kernel cheaper,
and \hh{} makes the whole group cheaper. A Hanner cube group contains all $2^\ell$ sign
vectors on its taps and finds the winner in $O(\ell)$ operations; a group of $k$ random
sign vectors needs $k$ dot products. Combining Hanner groups with NanoHydra's integer
pipeline is a natural next step. More broadly, ternary and binary networks
\citep{li2016ternary,rastegari2016xnor} and adder networks \citep{chen2020addernet} reduce
the cost of each learned weight; we do not learn weights and instead remove the dependence
on the number of kernels in a group.

\paragraph{Dictionary and pattern-based methods.} Dictionary methods classify series by
counting symbolic words, from SAX \citep{lin2007sax} to WEASEL~2.0
\citep{schaefer2023weasel2}, which is about as accurate as MiniRocket. Hanner features are
counts of words too: sign patterns on subsets of taps, weighted by the response. Sign
patterns of neighbouring samples are also the basis of one-dimensional local binary
patterns \citep{kaya2014lbp}. Ordinal patterns \citep{bandt2002permutation} are related but
different: they record the rank order of the samples in a window, not their signs.

\paragraph{Hashing and max-affine units.} The winner of a cube group is a sign pattern, the
code of hyperplane hashing \citep{charikar2002simhash}, and the winner of a cross-polytope
group is a signed coordinate, the code of cross-polytope hashing
\citep{andoni2015crosspolytope}; Section~\ref{sec:interpretation} makes this precise.
Winner-take-all hashing records the index of the largest of several coordinates
\citep{yagnik2011wta}, which is what Hydra's counts record for kernel responses. A group's
largest response, the support function of a polytope, is also the output of a maxout unit;
maxout polytopes \citep{balakin2025maxout} study the polytopes that networks of such units
define. We found no work that chooses competing kernels for the structure of their
polytope, or that finds a group's winner faster than by evaluating each kernel.

\paragraph{Evaluation.} We follow the bake-off protocol of 30 stratified resamples
\citep{ruiz2021multivariate,middlehurst2024bakeoff} on the UCR and UEA archives
\citep{dau2019ucr,bagnall2018uea}, with aeon's implementations and published results
\citep{middlehurst2024aeon}. Our primary test is the Bayesian signed-rank test of
\citet{benavoli2017bayesian}, which turns a margin of practical equivalence into a decision
rule, with Wilcoxon--Holm tests and critical-difference diagrams
\citep{demsar2006statistical} as secondary analyses; pairwise comparisons also avoid the
dependence of critical-difference diagrams on the set of compared methods
\citep{ismailfawaz2023mcm}. The UCR and UEA datasets are small, which favours
low-variance methods; the MONSTER collection \citep{dempster2025monster} offers large
datasets on which the speed of \hh{} would matter more.
